%!TEX root = main.tex
\documentclass[journal, onecolumn]{IEEEtran}
%\documentclass[12pt,journal,onecolumn]{IEEEtran}
%!TEX root = main.tex

% LC: debug
%==========================================================================
%\usepackage{refcheck}
%\usepackage[notref]{showkeys}
%==========================================================================

% LC: to be used for TRO
%==========================================================================
% \usepackage{mathptmx} % assumes new font selection scheme installed
%\usepackage{times} % assumes new font selection scheme installed
%==========================================================================
\usepackage{comment}
\usepackage[binary-units]{siunitx}
\usepackage{relsize}
\usepackage{ifthen}
\usepackage[colorinlistoftodos]{todonotes}

% \usepackage[caption=false]{subfig}

% \begin{comment}
% % Fancy formatting 
% \usepackage[tracking=false,kerning=true,spacing=true]{microtype}
% \usepackage[caption=false]{subfig}

% \usepackage[style=numeric-comp,sorting=none,firstinits=true, maxnames=3,bibstyle=numeric,abbreviate=false,defernums=true,eprint=false,backend=bibtex]{biblatex}
% \renewcommand{\bibfont}{\footnotesize}
% % \DeclareFieldFormat[article]{pages}{#1}%
% % \DeclareFieldFormat[inproceedings]{pages}{#1}%

% \AtEveryBibitem{%
%   \ifentrytype{article}{%
%     \clearfield{pages}%
%   }{%
%   }%
%   \ifentrytype{inproceedings}{%
%     \clearfield{pages}%
%   }{%
%   }%
% }
% \end{comment}

% \bibliography{commons/bib/strings_long}
% \bibliography{commons/bib/all_only_one_w_pages}
% \bibliography{../../references/references}

% \usepackage[noadjust]{cite}
\usepackage[vlined,ruled,linesnumbered]{algorithm2e}
\usepackage{graphics} % for pdf, bitmapped graphics files
\usepackage{rotating}
\usepackage{color}
\usepackage{enumerate}
\usepackage[T1]{fontenc}
\usepackage{psfrag}
\usepackage{epsfig} % for postscript graphics files
%\usepackage{subfigure}
% \usepackage{hyperref}
\usepackage{booktabs}
\usepackage{graphicx,url}
\usepackage{array}
\usepackage{latexsym}
\usepackage{amsfonts}
\usepackage{amsmath}
\usepackage{amssymb}
%\usepackage{amsthm}
\usepackage{xstring}
\usepackage[noend]{algorithmic}
\usepackage{multirow}
\usepackage{xcolor}
\usepackage{prettyref}
\usepackage{flexisym}
\usepackage{bigdelim}
\usepackage{breqn} % load this last
\usepackage{listings}
\let\labelindent\relax
\usepackage{enumitem}
\usepackage{xspace}
\usepackage{bm}
\graphicspath{{./figures/}}
\usepackage{tikz}
\usetikzlibrary{matrix,calc}


%\usepackage{ifpdf}
% Heiko Oberdiek's ifpdf.sty is very useful if you need conditional
% compilation based on whether the output is pdf or dvi.
% usage:
% \ifpdf
%   % pdf code
% \else
%    dvi code
% \fi
% The latest version of ifpdf.sty can be obtained from:
% http://www.ctan.org/tex-archive/macros/latex/contrib/oberdiek/
% Also, note that IEEEtran.cls V1.7 and later provides a builtin
% \ifCLASSINFOpdf conditional that works the same way.
% When switching from latex to pdflatex and vice-versa, the compiler may
% have to be run twice to clear warning/error messages.

% *** GRAPHICS RELATED PACKAGES ***
%
%\ifCLASSINFOpdf
  % \usepackage[pdftex]{graphicx}
  % declare the path(s) where your graphic files are
  % \graphicspath{{../pdf/}{../jpeg/}}
  % and their extensions so you won't have to specify these with
  % every instance of \includegraphics
  % \DeclareGraphicsExtensions{.pdf,.jpeg,.png}
%\else
  % or other class option (dvipsone, dvipdf, if not using dvips). graphicx
  % will default to the driver specified in the system graphics.cfg if no
  % driver is specified.
  % \usepackage[dvips]{graphicx}
  % declare the path(s) where your graphic files are
  % \graphicspath{{../eps/}}
  % and their extensions so you won't have to specify these with
  % every instance of \includegraphics
  % \DeclareGraphicsExtensions{.eps}
%\fi


\usepackage{mdwlist}
\let\stditemize\itemize
\let\endstditemize\enditemize
\let\itemize\undefined
\makecompactlist{itemize}{stditemize}

%\let\stdenumerate\enumerate
%\let\endstdenumerate\endenumerate
%\let\enumerate\undefined
%\makecompactlist{enumerate}{stdenumerate}
 

%!TEX root = main.tex

% LC: can be inserted
% \newcommand{\qed}{{\hfill $\square$}}

% Format definition
\newrefformat{prob}{Problem\,\ref{#1}}
\newrefformat{def}{Definition\,\ref{#1}}
\newrefformat{sec}{Section\,\ref{#1}}
\newrefformat{sub}{Section\,\ref{#1}}
\newrefformat{prop}{Proposition\,\ref{#1}}
\newrefformat{app}{Appendix\,\ref{#1}}
\newrefformat{alg}{Algorithm\,\ref{#1}}
\newrefformat{cor}{Corollary\,\ref{#1}}
\newrefformat{thm}{Theorem\,\ref{#1}}
\newrefformat{lem}{Lemma\,\ref{#1}}
\newrefformat{fig}{Fig.\,\ref{#1}}
\newrefformat{tab}{Table\,\ref{#1}}

% Problem environment
\newtheorem{theorem}{Theorem}
\newtheorem{problem}{Problem}
\newtheorem{trule}[theorem]{Rule}
\newtheorem{corollary}[theorem]{Corollary}
%\newtheorem{algorithm}{Algorithm}
%\newtheorem{procedure}{\textbf{Procedure}}
\newtheorem{conjecture}[theorem]{Conjecture}
\newtheorem{lemma}[theorem]{Lemma}
\newtheorem{assumption}[theorem]{Assumption}
\newtheorem{definition}[theorem]{Definition}
\newtheorem{proposition}[theorem]{Proposition}
\newtheorem{remark}[theorem]{Remark}
\newtheorem{example}[theorem]{Example}
\newtheorem{fact}[theorem]{Fact}

% Shortcuts
\newcommand{\cf}{\emph{cf.}\xspace}
\newcommand{\Figure}{Fig.}
\newcommand{\bdmath}{\begin{dmath}}
\newcommand{\edmath}{\end{dmath}}
\newcommand{\beq}{\begin{equation}}
\newcommand{\eeq}{\end{equation}}
\newcommand{\bdm}{\begin{displaymath}}
\newcommand{\edm}{\end{displaymath}}
\newcommand{\bea}{\begin{eqnarray}}
\newcommand{\eea}{\end{eqnarray}}
\newcommand{\beal}{\beq \begin{array}{ll}}
\newcommand{\eeal}{\end{array} \eeq}
\newcommand{\beas}{\begin{eqnarray*}}
\newcommand{\eeas}{\end{eqnarray*}}
\newcommand{\ba}{\begin{array}}
\newcommand{\ea}{\end{array}}
\newcommand{\bit}{\begin{itemize}}
\newcommand{\eit}{\end{itemize}}
\newcommand{\ben}{\begin{enumerate}}
\newcommand{\een}{\end{enumerate}}
\newcommand{\expl}[1]{&&\qquad\text{\color{gray}(#1)}\nonumber}

% \newcommand{\insertproof}[1]{%
%     % lem:name -> proof_lem_name.tex
%     \StrSubstitute{#1}{:}{_}[\name] % replace : with _
%     \begin{IEEEproof} %\color[rgb]{0.8,0.8,0.8}
%     \input{proof_\name.tex}
%     \end{IEEEproof}
% }

% Calligraphic fonts
\newcommand{\calA}{{\cal A}}
\newcommand{\calB}{{\cal B}}
\newcommand{\calC}{{\cal C}}
\newcommand{\calD}{{\cal D}}
\newcommand{\calE}{{\cal E}}
\newcommand{\calF}{{\cal F}}
\newcommand{\calG}{{\cal G}}
\newcommand{\calH}{{\cal H}}
\newcommand{\calI}{{\cal I}}
\newcommand{\calJ}{{\cal J}}
\newcommand{\calK}{{\cal K}}
\newcommand{\calL}{{\cal L}}
\newcommand{\calM}{{\cal M}}
\newcommand{\calN}{{\cal N}}
\newcommand{\calO}{{\cal O}}
\newcommand{\calP}{{\cal P}}
\newcommand{\calQ}{{\cal Q}}
\newcommand{\calR}{{\cal R}}
\newcommand{\calS}{{\cal S}}
\newcommand{\calT}{{\cal T}}
\newcommand{\calU}{{\cal U}}
\newcommand{\calV}{{\cal V}}
\newcommand{\calW}{{\cal W}}
\newcommand{\calX}{{\cal X}}
\newcommand{\calY}{{\cal Y}}
\newcommand{\calZ}{{\cal Z}}

% SETS:
\newcommand{\setA}{\textsf{A}}
\newcommand{\setB}{\textsf{B}}
\newcommand{\setC}{\textsf{C}}
\newcommand{\setD}{\textsf{D}}
\newcommand{\setE}{\textsf{E}}
\newcommand{\setF}{\textsf{F}}
\newcommand{\setG}{\textsf{G}}
\newcommand{\setH}{\textsf{H}}
\newcommand{\setI}{\textsf{I}}
\newcommand{\setJ}{\textsf{J}}
\newcommand{\setK}{\textsf{K}}
\newcommand{\setL}{\textsf{L}}
\newcommand{\setM}{\textsf{M}}
\newcommand{\setN}{\textsf{N}}
\newcommand{\setO}{\textsf{O}}
\newcommand{\setP}{\textsf{P}}
\newcommand{\setQ}{\textsf{Q}}
\newcommand{\setR}{\textsf{R}}
\newcommand{\setS}{\textsf{S}}
\newcommand{\setT}{\textsf{T}}
\newcommand{\setU}{\textsf{U}}
\newcommand{\setV}{\textsf{V}}
\newcommand{\setX}{\textsf{X}}
\newcommand{\setY}{\textsf{Y}}
\newcommand{\setZ}{\textsf{Z}}

%General
\newcommand{\smallheading}[1]{\textit{#1}: }
\newcommand{\algostep}[1]{{\small\texttt{#1:}}\xspace}
\newcommand{\etal}{\emph{et~al.}\xspace}
\newcommand{\setal}{~\emph{et~al.}\xspace}
\newcommand{\eg}{\emph{e.g.,}\xspace}
\newcommand{\ie}{\emph{i.e.,}\xspace}
\newcommand{\myParagraph}[1]{{\bf #1.}\xspace}
% \newcommand{\email}[1]{{\smaller \textsf{#1}}}

%Typography
\newcommand{\M}[1]{{\bm #1}} % Face for matrices
\renewcommand{\boldsymbol}[1]{{\bm #1}}
\newcommand{\algoname}[1]{\textsc{#1}} % Name of algorithms

%Editing
\newcommand{\towrite}{{\color{red} To write}\xspace}
\newcommand{\xxx}{{\color{red} XXX}\xspace}
\newcommand{\AC}[1]{{\color{red} \textbf{AC}: #1}}
\newcommand{\LC}[1]{{\color{red} \textbf{LC}: #1}}
\newcommand{\FM}[1]{{\color{red} \textbf{FM}: #1}}
\newcommand{\RT}[1]{{\color{blue} \textbf{RT}: #1}}
\newcommand{\VT}[1]{{\color{blue} \textbf{VT}: #1}}
\newcommand{\hide}[1]{}
\newcommand{\wrt}{w.r.t.\xspace}
\newcommand{\highlight}[1]{{\color{red} #1}}
\newcommand{\tocheck}[1]{{\color{brown} #1}}
\newcommand{\grayout}[1]{{\color{gray} #1}}
\newcommand{\grayText}[1]{{\color{gray} \text{#1} }}
\newcommand{\grayMath}[1]{{\color{gray} #1 }}
\newcommand{\hiddenText}{{\color{gray} hidden text.}}
\newcommand{\hideWithText}[1]{\hiddenText}
\newcommand{\marginNote}[1]{\marginpar{\small \red{#1}}}


\newcommand{\NA}{{\sf n/a}}
\newcommand{\versus}{\scenario{VS}\xspace}

%Basic math symbols
\newcommand{\kron}{\otimes}
\newcommand{\dist}{\mathbf{dist}}
\newcommand{\iter}{\! \rm{iter.} \;}
\newcommand{\leqt}{\!\!\! < \!\!\!}
\newcommand{\geqt}{\!\!\! > \!\!\!}
\newcommand{\mysetminus}{-} % One set minus another
\newcommand{\powerset}{\mathcal{P}}
\newcommand{\Int}[1]{ { {\mathbb Z}^{#1} } }
\newcommand{\Natural}[1]{ { {\mathbb N}^{#1} } }
\newcommand{\Complex}[1]{ { {\mathbb C}^{#1} } }
\newcommand{\one}{ {\mathbf{1}} }
\newcommand{\subject}{\text{ subject to }}
\DeclareMathOperator*{\argmax}{arg\,max}
\DeclareMathOperator*{\argmin}{arg\,min}

%% Norms
\newcommand{\normsq}[2]{\left\|#1\right\|^2_{#2}}
\newcommand{\norm}[1]{\left\| #1 \right\|}
\newcommand{\normsqs}[2]{\|#1\|^2_{#2}}
\newcommand{\infnorm}[1]{\|#1\|_{\infty}}
\newcommand{\zeronorm}[1]{\|#1\|_{0}}
\newcommand{\onenorm}[1]{\|#1\|_{1}}
\newcommand{\lzero}{\ell_{0}}
\newcommand{\lone}{\ell_{1}}
\newcommand{\linf}{\ell_{\infty}}

\newcommand{\E}{{\mathbb{E}}}
\newcommand{\EV}{\mathbb{E}}
\newcommand{\erf}{{\mathbf{erf}}}
\newcommand{\prob}[1]{{\mathbb P}\left(#1\right)}
% \newcommand{\tran}{^{\top}}
\newcommand{\tran}{^{\mathsf{T}}}
\newcommand{\traninv}{^{-\mathsf{T}}}
\newcommand{\diag}[1]{\mathrm{diag}\left(#1\right)}
\newcommand{\trace}[1]{\mathrm{tr}\left(#1\right)}
\newcommand{\conv}[1]{\mathrm{conv}\left(#1\right)}
\newcommand{\polar}[1]{\mathrm{polar}\left(#1\right)}
\newcommand{\rank}[1]{\mathrm{rank}\left(#1\right)}
\newcommand{\e}{{\mathrm e}}
\newcommand{\inv}{^{-1}}
\newcommand{\pinv}{^\dag}
\newcommand{\until}[1]{\{1,\dots, #1\}}
\newcommand{\ones}{{\mathbf 1}}
\newcommand{\zero}{{\mathbf 0}}
\newcommand{\eye}{{\mathbf I}}
\newcommand{\vect}[1]{\left[\begin{array}{c}  #1  \end{array}\right]}
\newcommand{\matTwo}[1]{\left[\begin{array}{cc}  #1  \end{array}\right]}
\newcommand{\matThree}[1]{\left[\begin{array}{ccc}  #1  \end{array}\right]}
\newcommand{\dss}{\displaystyle}
\newcommand{\Real}[1]{ { {\mathbb R}^{#1} } }
\newcommand{\reals}{\Real{}}
\newcommand{\opt}{^{\star}}
\newcommand{\only}{^{\alpha}}
\newcommand{\copt}{^{\text{c}\star}}
\newcommand{\atk}{^{(k)}}
\newcommand{\att}{^{(t)}}
\newcommand{\at}[1]{^{(#1)}}
\newcommand{\attau}{^{(\tau)}}
\newcommand{\atc}[1]{^{(#1)}}
\newcommand{\atK}{^{(K)}}
\newcommand{\atj}{^{(j)}}
\newcommand{\projector}{{\tt projector}}
\newcommand{\setdef}[2]{ \{#1 \; {:} \; #2 \} }
\newcommand{\smalleye}{\left(\begin{smallmatrix}1&0\\0&1\end{smallmatrix}\right)}

%Spaces
\newcommand{\SEtwo}{\ensuremath{\mathrm{SE}(2)}\xspace}
\newcommand{\SE}[1]{\ensuremath{\mathrm{SE}(#1)}\xspace}
\newcommand{\SEthree}{\ensuremath{\mathrm{SE}(3)}\xspace}
\newcommand{\SOtwo}{\ensuremath{\mathrm{SO}(2)}\xspace}
\newcommand{\SOthree}{\ensuremath{\mathrm{SO}(3)}\xspace}
\newcommand{\Othree}{\ensuremath{\mathrm{O}(3)}\xspace}
\newcommand{\SOn}{\ensuremath{\mathrm{SO}(n)}\xspace}
\newcommand{\SO}[1]{\ensuremath{\mathrm{SO}(#1)}\xspace}
\newcommand{\On}{\ensuremath{\mathrm{O}(n)}\xspace}
\newcommand{\sotwo}{\ensuremath{\mathrm{so}(2)}\xspace}
\newcommand{\sothree}{\ensuremath{\mathrm{so}(3)}\xspace}
\newcommand{\intexpmap}[1]{\mathrm{Exp}\left(#1\right)}
\newcommand{\intlogmap}[1]{\mathrm{Log}\left(#1\right)}
\newcommand{\logmapz}[1]{\mathrm{Log}_0(#1)}
\newcommand{\loglikelihood}{\mathrm{log}\calL}
\newcommand{\intprinlogmap}{\mathrm{Log}_0}
\newcommand{\expmap}[1]{\intexpmap{#1}}
\newcommand{\expmaps}[1]{\langle #1 \rangle_{2\pi}}
\newcommand{\biggexpmap}[1]{\left\langle #1 \right\rangle_{2\pi}}
\newcommand{\logmap}[1]{\intlogmap{#1}}
\newcommand{\round}[1]{\mathrm{round}\left( #1 \right)}

% Matrices 
\newcommand{\MA}{\M{A}}
\newcommand{\MB}{\M{B}}
\newcommand{\MC}{\M{C}}
\newcommand{\MD}{\M{D}}
\newcommand{\ME}{\M{E}}
\newcommand{\MJ}{\M{J}}
\newcommand{\MK}{\M{K}}
\newcommand{\MG}{\M{G}}
\newcommand{\MM}{\M{M}}
\newcommand{\MN}{\M{N}}
\newcommand{\MP}{\M{P}}
\newcommand{\MQ}{\M{Q}}
\newcommand{\MU}{\M{U}}
\newcommand{\MR}{\M{R}}
\newcommand{\MS}{\M{S}}
\newcommand{\MI}{\M{I}}
\newcommand{\MV}{\M{V}}
\newcommand{\MF}{\M{F}}
\newcommand{\MH}{\M{H}}
\newcommand{\ML}{\M{L}}
\newcommand{\MO}{\M{O}}
\newcommand{\MT}{\M{T}}
\newcommand{\MX}{\M{X}}
\newcommand{\MY}{\M{Y}}
\newcommand{\MW}{\M{W}}
\newcommand{\MZ}{\M{Z}}
\newcommand{\MSigma}{\M{\Sigma}} 
\newcommand{\MOmega}{\M{\Omega}}
\newcommand{\MPhi}{\M{\Phi}}
\newcommand{\MPsi}{\M{\Psi}}
\newcommand{\MDelta}{\M{\Delta}}
\newcommand{\MLambda}{\M{\Lambda}}

% vectors
\newcommand{\va}{\boldsymbol{a}} 
\newcommand{\vh}{\boldsymbol{h}} 
\newcommand{\vb}{\boldsymbol{b}}
\newcommand{\vc}{\boldsymbol{c}}
\newcommand{\vd}{\boldsymbol{d}}
\newcommand{\ve}{\boldsymbol{e}}
\newcommand{\vf}{\boldsymbol{f}}
\newcommand{\vg}{\boldsymbol{g}}
\newcommand{\vk}{\boldsymbol{k}}
\newcommand{\vl}{\boldsymbol{l}}
\newcommand{\vn}{\boldsymbol{n}}
\newcommand{\vo}{\boldsymbol{o}}
\newcommand{\vp}{\boldsymbol{p}}
\newcommand{\vq}{\boldsymbol{q}}
\newcommand{\vr}{\boldsymbol{r}}
\newcommand{\vs}{\boldsymbol{s}}
\newcommand{\vu}{\boldsymbol{u}}
\newcommand{\vv}{\boldsymbol{v}}
\newcommand{\vt}{\boldsymbol{t}}
\newcommand{\vxx}{\boldsymbol{x}} 
\newcommand{\vy}{\boldsymbol{y}}
\newcommand{\vw}{\boldsymbol{w}}
\newcommand{\vzz}{\boldsymbol{z}} 
\newcommand{\vdelta}{\boldsymbol{\delta}}
\newcommand{\vgamma}{\boldsymbol{\gamma}}  
\newcommand{\vlambda}{\boldsymbol{\lambda}}
\newcommand{\vtheta}{\boldsymbol{\theta}}
\newcommand{\valpha}{\boldsymbol{\alpha}}
\newcommand{\vbeta}{\boldsymbol{\beta}}
\newcommand{\vnu}{\boldsymbol{\nu}}
\newcommand{\vmu}{\boldsymbol{\mu}}
\newcommand{\vepsilon}{\boldsymbol{\epsilon}}
\newcommand{\vtau}{\boldsymbol{\tau}}

%Intrinsic geometry
\newcommand{\Rtheta}{\boldsymbol{R}}
\newcommand{\symf}{f} % Symmetry function

%Angles
\newcommand{\angledomain}{(-\pi,+\pi]}

% Tree, graphs, and cycle basis
\newcommand{\MCB}{\mathsf{MCB}}
\newcommand{\FCM}{\mathsf{FCM}}
\newcommand{\FCB}{\mathsf{FCB}}
\newcommand{\cyclemap}[1]{\calC^{\calG}\left(#1\right)}
\newcommand{\incidencemap}[1]{\calA^{\calG}\left(#1\right)}
\newcommand{\cyclemapk}{\calC^{\calG}_{k}}
\newcommand{\incidencemapij}{\calA^{\calG}_{ij}}
\renewcommand{\ij}{_{ij}}
\newcommand{\foralledges}{\forall(i,j) \in \calE}
\newcommand{\sumalledges}{
     \displaystyle
     \sum_{(i,j) \in \calE}}
\newcommand{\sumalledgesm}{
     \displaystyle
     \sum_{i=1}^{m}}
\newcommand{\T}{\mathsf{T}}
\newcommand{\To}{\T_{\rm o}}
\newcommand{\Tm}{\T_{\rm m}}
\newcommand{\MCBa}{\MCB_{\mathsf{a}}}
\newcommand{\FCBo}{\FCB_{\mathsf{o}}}
\newcommand{\FCBm}{\FCB_{\mathsf{m}}}

% Algorithms
\newcommand{\algoonlyname}{MOLE2D}
\newcommand{\algoml}{{\smaller\sf \algoonlyname}\xspace}
\newcommand{\algocyclebasis}{\algoname{compute-cycle-basis}}
\newcommand{\scenario}[1]{{\smaller \sf#1}\xspace}
\newcommand{\toro}{{\smaller\sf Toro}\xspace}
\newcommand{\gtwoo}{{\smaller\sf g2o}\xspace}
\newcommand{\gtwooST}{{\smaller\sf g2oST}\xspace}
\newcommand{\gtwood}{{\smaller\sf g2o{10}}\xspace}
\newcommand{\gtsam}{{\smaller\sf gtsam}\xspace}
\newcommand{\isam}{{\smaller\sf iSAM}\xspace}
\newcommand{\lago}{{\smaller\sf LAGO}\xspace}
\newcommand{\egtwoo}{{\smaller\sf \algoonlyname+g2o}\xspace}

% Datasets
%\newcommand{\grid}{\scenario{cube}}
\newcommand{\rim}{\scenario{rim}}
\newcommand{\cubicle}{\scenario{cubicle}}
\newcommand{\sphere}{\scenario{sphere}}
\newcommand{\sphereHard}{\scenario{sphere-a}}
\newcommand{\garage}{\scenario{garage}}
\newcommand{\torus}{\scenario{torus}}
\newcommand{\oneloop}{\scenario{circle}}
\newcommand{\intel}{\scenario{INTEL}}
\newcommand{\bovisa}{\scenario{Bovisa}}
\newcommand{\bov}{\scenario{B25b}}
\newcommand{\fra}{\scenario{FR079}}
\newcommand{\frb}{\scenario{FRH}}
\newcommand{\csail}{\scenario{CSAIL}}
\newcommand{\Ma}{\scenario{M3500}}
\newcommand{\Mb}{\scenario{M10000}}
\newcommand{\ATE}{\scenario{ATE}}
\newcommand{\CVX}{\scenario{CVX}}
\newcommand{\NEOS}{\scenario{NEOS}}
\newcommand{\sdptThree}{\scenario{sdpt3}}
\newcommand{\MOSEK}{\scenario{MOSEK}}
\newcommand{\NESTA}{\scenario{NESTA}}
\newcommand{\vertigo}{\scenario{Vertigo}}
\newcommand{\SDPA}{\scenario{SDPA}}

% \newcommand{\tablabel}[1]{%
%     \hspace{-1mm}%
%     \begin{sideways}{\small\scenario{#1}}\end{sideways}%
%     \hspace{-4mm}%
% }
%
% \newcommand{\tabfig}[2]{%
%     \subfloat[\label{fig:#2} #1]{%
%     \begin{minipage}{5.5cm}%
%     \centering%
%     \includegraphics[width=5.5cm,trim=4 0 4 15]{figures/#2}
%     \end{minipage}%
%     }%
% }
%
% \newcommand{\subFigure}[3]{%
%     \subfloat[\label{fig:#2} #1]{%
%     \begin{minipage}{#3cm}%
%     \centering%
%     \includegraphics[width=#3cm,trim=4 0 4 15]{figures/#2}
%     \end{minipage}%
%     }%
% }
\newcommand{\cvx}{{\sf cvx}\xspace}

% COLORS
\newcommand{\blue}[1]{{\color{blue}#1}}
\newcommand{\green}[1]{{\color{green}#1}}
\newcommand{\red}[1]{{\color{red}#1}}

% TO MANAGE REFERENCES
%============================================================================
\newcommand{\linkToPdf}[1]{\href{#1}{\blue{(pdf)}}}
\newcommand{\linkToPpt}[1]{\href{#1}{\blue{(ppt)}}}
\newcommand{\linkToCode}[1]{\href{#1}{\blue{(code)}}}
\newcommand{\linkToWeb}[1]{\href{#1}{\blue{(web)}}}
\newcommand{\linkToVideo}[1]{\href{#1}{\blue{(video)}}}
\newcommand{\linkToMedia}[1]{\href{#1}{\blue{(media)}}}
\newcommand{\award}[1]{\xspace} % {{\red{#1}}} % omit awards


% PAPER-SPECIFIC COMMANDS
%============================================================================

% Linear approximation
\newcommand{\vpose}{\boldsymbol{x}}
\newcommand{\vz}{\boldsymbol{z}}
\newcommand{\vDelta}{\boldsymbol{\Delta}}
\newcommand{\vposesub}{\hat{\vpose}}
\newcommand{\vpossub}{\hat{\vpos}}
\newcommand{\vthetasub}{\hat{\vtheta}}
\newcommand{\Pthetasub}{\MP_\vtheta}
\newcommand{\thetasub}{\hat{\theta}}
\newcommand{\vposecorr}{\tilde{\vpose}}
\newcommand{\vposcorr}{\tilde{\vpos}}
\newcommand{\vthetacorr}{\tilde{\vtheta}}
\newcommand{\vposestar}{{\vpose}^{\star}}
\newcommand{\vposstar}{{\vpos}^{\star}}
\newcommand{\vthetastar}{{\vtheta}^{\star}}
\newcommand{\vcthetastar}{{\vctheta}^{\star}}
\newcommand{\pose}{\boldsymbol{x}}
\newcommand{\pos}{\boldsymbol{p}}
\newcommand{\mease}{z} % element
\newcommand{\meas}{\boldsymbol{z}} % vector
\newcommand{\meashat}{\hat{\meas}}




%!TEX root = main.tex
\lstdefinelanguage{mine} 
{morekeywords={while,True,if,break,=,return,function,for,until,in,input,output,assumptions,assumption,invariant,loop,variant,end,invariant,precondition,variables}, 
sensitive=false, 
morecomment=[l]{\#}, 
morecomment=[il]\%{.}, 
morecomment=[s]{/*}{*/}, 
morestring=[b]", 
} 

\lstset{% general command to set parameter(s) 
basicstyle=\small\rmfamily, % print whole listing small 
keywordstyle=\color{black}\bfseries, 
% underlined bold black keywords 
identifierstyle=, % nothing happens 
commentstyle=\color[rgb]{0.2,0.2,0.2},
stringstyle=\ttfamily, % typewriter type for strings 
showstringspaces=false,
mathescape=true,
fontadjust=true,
columns=fullflexible} % no special string spaces 
\lstset{frameround=fttt,frame=none}

\usepackage{calc}
\newlength\listingnumberwidth
\setlength\listingnumberwidth{\widthof{00}}
\lstset{
numbers=left,
xleftmargin=\listingnumberwidth,
numbersep=1em
}


\lstset{escapeinside={*@}{@*}}

\newcommand{\tuple}[2]{\langle#1,#2\rangle}


\usepackage{comment}
\usepackage{cite}

\newcommand{\point}{\vp}
\newcommand{\points}{\vp}

%\renewcommand\baselinestretch{0.96}

\begin{document}

 \title{\LARGE{Scheduling the RSS 2025 Area Chair Meeting \\ via Mixed-Integer Linear Programming}}

\author{Luca Carlone and Lukas Schmid
\thanks{L.\,Carlone is with the Laboratory for Information \& Decision Systems, Massachusetts Institute of Technology, Cambridge, MA, USA, 
{\sf lcarlone@mit.edu}
}}

\newcommand{\MRij}{\bar\MR\ij}
\newcommand{\MRijy}{\bar\MR_{ij,y}}
\newcommand{\MRiy}{\MR_{y}(\theta_i) }
\newcommand{\MRip}{\bar\MR_{i,p}}
\newcommand{\MRir}{\bar\MR_{i,r}}
\newcommand{\MRjy}{\MR_{y}(\theta_j) }
\newcommand{\MRjp}{\bar\MR_{j,p}}
\newcommand{\MRjr}{\bar\MR_{j,r}}
\newcommand{\vtij}{\bar\vt\ij}
\newcommand{\vtijy}{\bar\vt_{ij,y}}
\newcommand{\Log}[1]{\text{Log}\left(#1\right)}

\maketitle

In this document we solve the problem of scheduling the sessions of the RSS 2025 Area Chair meeting as an Mixed-Integer Linear Program.
In particular, we simultaneously compute assignments of papers and ACs to discussion sessions.

\section{General Formulation}

\myParagraph{Variables}
The decision variables in this problem are the \emph{AC-to-session} assignment matrix $\MX$ and 
the \emph{paper-to-session} assignment matrix $\MY$.  
In particular, $\MX \in \{0,1\}^{n \times s}$, where $n$ is the number of Area Chairs (ACs) and $s$ is the number of sessions, and is such that $\MX_{ij} = 1$ if AC $i$ is assigned to session $j$.
In our case, $n = 54$ and $s = 36$ (4 paper discussion slots in the schedule, each with 9 parallel breakouts, 8 for discussion + 1 ``standout'' breakout for ACs that cannot join any of the regular sessions).
%
The matrix $\MY \in \{0,1\}^{p \times s}$, where $p$ is the number of papers to discuss, is such that $\MY_{kj}$ if paper $k$ is assigned to session $j$.
In our case, $s = 168$ (\ie we expect to discuss 6 papers in each of the 3 main discussion sessions, plus 3 in each overflow session).  
Our problem below also computes an auxiliary variable $\MZ \in \Real{n \times s}$, that is such that $\MZ_{ij}$ is equal to the cumulative affinity score (sum of the affinity scores for all papers in that session) of AC $i$ and session $j$ if the AC is assigned to that session, or is zero otherwise.
% if $i$ is not assigned to $j$.

\myParagraph{Inputs}
We are given an \emph{affinity} matrix $\MA \in \Real{n \times p}$, which contains the affinity of the ACs for all papers to discuss. In other words, 
$\MA_{ik}$ is the affinity of AC $i$ with paper $k$. 
Moreover, we are given a ``primary AC assignment'' matrix  $\MP \in \Real{n \times p}$, that is such that $\MP_{ik} = 1$ if AC $i$ is the primary handling AC for paper $k$.
 Similarly, we are given a ``secondary AC assignment'' matrix  $\MS \in \Real{n \times p}$, that is such that $\MS_{ik} = 1$ if AC $i$ is the secondary handling AC for paper $k$.
Finally, we are given a Conflict of Interest (COI) matrix $\MC \in \Real{n \times p}$, that is such that $\MC_{ik} = 1$ if AC $i$ is has a COI with paper $k$.
 
\myParagraph{Objective} Therefore, the goal is to simultaneously compute paper-to-session assignments and AC-to-session assignments such that we maximize the affinity of ACs between papers and ACs in a session, while avoiding COIs and ensuring that when a paper is assigned to a session, then its primary and secondary ACs are in that session as well.

\myParagraph{Formulation} The optimization problem computes the optimal matching can be states as follows:

\begin{align}
\max_{\substack{\MX \in \{0,1\}^{n \times s} \\  \MY \in \{0,1\}^{p \times s} \\  \MZ \in \real{n \times s}} } & \ones\tran \MZ \ones & \text{(maximize cumulative affinity score)} \nonumber \\
\subject &  \MX \cdot \ve_{[1:9]} = \ones_n  & \text{(each AC join only 1 breakout for the 1st discussion slot)}\nonumber  \\
         &  \MX \cdot \ve_{[10:18]} = \ones_n  & \text{(each AC join only 1 breakout for the 2nd discussion slot)} \nonumber \\
         &  \MX \cdot \ve_{[19:27]} = \ones_n  & \text{(each AC join only 1 breakout for the 3rd discussion slot)} \nonumber \\
         &  \MX \cdot \ve_{[28:36]} = \ones_n  & \text{(each AC join only 1 breakout for the 4th discussion slot)} \nonumber \\
         &  (\ones\tran \cdot \MX \leq \vn_\max\tran) & \text{(optional: at most $\vn_\max$ ACs in a breakout)} \nonumber \\
         &  (\ones\tran \cdot \MX \geq \vn_\min\tran) & \text{(optional: at least $\vn_\min$ ACs in a breakout)} \nonumber \\
         \nonumber \\
         &  \ones\tran \cdot \MY \leq \vp_{max}\tran  & \text{(at most 6 papers per regular session, 4 in the overflow sessions, 0 in standout)} \nonumber \\
         &  \MY \cdot \ones  = \ones  & \text{(each paper is assigned to a unique session)} \nonumber \\
         % &  \MY \cdot \ve_{9} = \zero  & \text{(each paper is only assigned to 1 session)} \\
         % &  \MY \cdot \ve_{18} = \zero  & \text{(papers cannot be assigned to overflow session)} \\
         % &  \MY \cdot \ve_{27} = \zero  & \text{(each paper is only assigned to 1 session)} \\
         % &  \MY \cdot \ve_{36} = \zero  & \text{(papers cannot be assigned to overflow session)} \\
         \nonumber \\
         &  \MX \leq \ones\tran \ones - \frac{1}{p} (\MC \cdot \MY)  & \text{(avoid COIs)} \nonumber \\
         &  \MX \geq \frac{1}{p} (\MP \cdot \MY)  & \text{(primary AC must be in session where paper is discussed)} \nonumber \\
         &  \MX \geq \frac{1}{p} (\MS \cdot \MY)  & \text{(secondary AC must be in session where paper is discussed)} \nonumber \\
         \nonumber \\ 
         & \MZ \leq a_\max \MX & \text{(force affinities to zero whenever an AC $i$ is not in session $j$)} \nonumber \\
         & \MZ \leq \MA \MY & \text{(forces an upper-bound on the affinities, which is attained when AC $i$ is in session $j$)} \nonumber 
\end{align}
where $\ones_n$ is the vector of length $n$ with all entries equal to one, 
$\vp_{max} \in \Real{s}$ is a vector whose $j$-th entry is the maximum number of papers that can be discussed in breakout $j$ (this is zero for the standout breakout),
$\vn_\max \in \Real{s}$ is a vector whose $j$-th entry is the maximum number of AC that can be in breakout $j$ (this is $n$ for the standout breakout),
$\vn_\min \in \Real{s}$ is a vector whose $j$-th entry is the minimum number of AC that can be in breakout $j$ (this is $0$ for the standout breakout),
$\ve_{[\vv]}$ is the vector with entries equal to $1$ for the indices in $\vv$ and zero otherwise, 
and $a_\max$ is a constant (the maximum cumulative affinity score).

Let us explain the optimization in more detail.
The matrix $\MZ$ represents the cumulative affinity scores for each AC to all papers in a session (and is nonzero only for the sessions the AC is assigned to).
Therefore, the objective $\ones\tran \MZ \ones$ simply sums and maximizes the cumulative affinity scores across all assignments.
Many of the constraints are self-explanatory from the description above, so we focus on the most complex ones.
  The constraint $\MX \leq \ones\tran \ones - \frac{1}{p} (\MC \cdot \MY)$ enforces that ACs are not assigned to sessions for which they have at least a COI.
  In particular the matrix $\MC \cdot \MY$ is such that, the entry $i,j$ of the matrix, namely $[\MC \cdot \MY]_{ij}$, describes the number of COI AC $i$ has with session $j$ (this is at most $p$, and typically much smaller than that).
  Therefore, whenever $[\frac{1}{p} (\MC \cdot \MY)]_{ij} > 0$, then $[\ones\tran \ones - \frac{1}{p} (\MC \cdot \MY)]_{ij} < 1$ and the corresponding entry of $\MX$ is forced to zero.
For constraint $\MX \geq \frac{1}{p} (\MP \cdot \MY)$, the matrix $\MP \cdot \MY$ is such that, the entry $i,j$ of the matrix, namely $[\MP \cdot \MY]_{ij}$, describes the number of papers AC $i$ is the primary AC for in session $j$ (this is at most $p$, and typically much smaller than that). 
Therefore, whenever $[\frac{1}{p} (\MC \cdot \MY)]_{ij} > 0$, the corresponding entry of $\MX$ is forced to be one (meaning that the AC is forced to join that session).
Together, the last two constraints are the same as $\MZ \leq \min(a_\max \MX,\MA \MY)$ (where the minimum is computed entry-wise), which forces $\MZ_{ij}=0$   whenever $\MX_{ij}=0$ (\ie we don't count affinities when an AC is not assigned to a session), and $\MZ_{ij} = [\MA \MY]_{ij}$ when $\MX_{ij}=1$ (recall that $a_\max \gg [\MA \MY]_{ij}$ for any $i$ and $j$). 

In Matlab's {\tt intlinprog}, the decision variable must be a vector, in our case $\vxx = [\text{vec}(\MX) \; \text{vec}(\MY) \; \text{vec}(\MZ)]$, hence we rewrite the previous problem as:

\newcommand{\myQuad}{\qquad}
\begin{align}
\max_{\vxx \in \{0,1\}^{2 ns + ps}} & \ones_{ns}\tran \cdot \ME_\MZ \cdot \vxx  \myQuad \text{(maximize cumulative affinity score)} \nonumber \\
\subject &  (\ve_{[1:9]}\tran \kron \eye_n) \cdot \ME_\MX \cdot \vxx = \ones_n  \myQuad \text{(each AC join only 1 breakout for the 1st discussion slot)}\nonumber  \\
         &  (\ve_{[10:18]}\tran \kron \eye_n) \cdot \ME_\MX \cdot \vxx = \ones_n   \myQuad \text{(each AC join only 1 breakout for the 2nd discussion slot)} \nonumber \\
         &  (\ve_{[19:27]}\tran \kron \eye_n) \cdot \ME_\MX \cdot \vxx = \ones_n  \myQuad \text{(each AC join only 1 breakout for the 3rd discussion slot)} \nonumber \\
         & (\ve_{[28:36]}\tran \kron \eye_n) \cdot \ME_\MX \cdot \vxx = \ones_n  \myQuad \text{(each AC join only 1 breakout for the 4th discussion slot)} \nonumber \\
         &  (\eye_s \kron \ones_n\tran)  \cdot \ME_{\MX} \cdot \vxx \leq \vn_\max) \myQuad \text{(optional: at most $\vn_\max$ ACs in a breakout)} \nonumber \\
         &  (\eye_s \kron \ones_n\tran) \cdot \ME_{\MX} \cdot \vxx  \geq \vn_\min) \myQuad \text{(optional: at least $\vn_\min$ ACs in a breakout)} \nonumber \\
         \nonumber \\
         & (\eye_s \kron \ones_p\tran ) \cdot \ME_{\MY} \cdot \vxx \leq \vp_{max}  \myQuad \text{(at most 6 papers per session, 4 in the overflow sessions, 0 in standout)} \nonumber \\
         & (\ones_{s}\tran \kron \eye_p) \cdot \ME_{\MY} \cdot \vxx  = \ones_p  \myQuad \text{(each paper is assigned to a unique session)} \nonumber \\
         \nonumber \\
         &  \ME_\MX \cdot \vxx \leq \ones_{ns} - \frac{1}{p} (\eye_s \kron \MC) \cdot \ME_{\MY} \cdot \vxx  \myQuad \text{(avoid COIs)} \nonumber \\
         &  \ME_\MX \cdot \vxx \geq \frac{1}{p} (\eye_s \kron \MP) \cdot \ME_{\MY} \cdot  \vxx  \myQuad \text{(primary AC must be in session where paper is discussed)} \nonumber \\
         &  \ME_\MX \cdot \vxx \geq \frac{1}{p} (\eye_s \kron \MS) \cdot \ME_{\MY} \cdot \vxx \myQuad \text{(secondary AC must be in session where paper is discussed)} \nonumber \\
         \nonumber \\ 
         & \ME_\MZ \cdot \vxx  \leq a_\max \cdot \ME_\MX \cdot \vxx \myQuad \text{(force affinities to zero whenever an AC $i$ is not in session $j$)} \nonumber \\
         & \ME_\MZ \cdot \vxx \leq (\eye_s \kron \MA) \cdot \ME_\MY \cdot \vxx  \myQuad \text{(forces an upper-bound on the affinities, which is attained when AC $i$ is in session $j$)} \nonumber 
\end{align}
where $\va = \text{vec}(\MA)$ (column-wise vectorization), $\eye_n$ is the identity matrix of size $n$, and we used the property 
that for two matrices $\MF \in \Real{a \times b}$ and $\MG \in \Real{b \times c}$ it holds $\text{vec}(\MF \cdot \MG) =  (\MG\tran \kron \eye_a) \cdot \text{vec}(\MF) = 
(\eye_c \kron \MF) \cdot \text{vec}(\MG)$. 
Moreover, we defined the selection matrices $\ME_\MX = [\eye_{ns} \;\; \zero_{ns,ps} \;\; \zero_{ns,ns}]$, $\ME_\MY  = [\zero_{ps,ns} \;\; \eye_{ps} \;\; \zero_{ps,ns}]$, and   $\ME_\MZ = [\zero_{ns,ns} \;\; \zero_{ns,ps} \;\; \eye_{ns}]$ that select the entries corresponding to $\MX$, $\MY$, and $\MZ$ from $\vxx$, respectively.

\myParagraph{Adding Senior Area Chairs} Senior Area Chairs have to satisfy similar constraints as the ACs, but have to satisfy the additional constraint that 
each AC has to staff a different breakout and none of them should be in a standout breakout.
Assuming the Senior AC are listed as the last $n_{\text{sac}}$ rows of $\MX$, these constraints can be written as follows:
\begin{align}
\ve_{[\text{end}-n_{\text{sac}}:\text{end}]}\tran \cdot \MX = \vs_{\text{sac}}\tran
\end{align}
where $\vs_{\text{sac}}$ is $1$ for the regular sessions and zero for the standout sessions.
After vectorization, this constraint becomes:
\begin{align}
(\eye_s \kron \ve_{[\text{end}-n_{\text{sac}}:\text{end}]}\tran) \cdot \ME_\MX \cdot \vxx = \vs_{\text{sac}}
\end{align}

% \myParagraph{Maximizing the Minimum Affinity instead of the Cumulative Affinity} %% this is currently incorrect, since the extra variable is always zero
% In this case, we want to maximize the minimum of $\MZ$ instead of the sum of its entries.
% We can do so, by adding an extra variable at the end of $\vxx$, and then maximize it in the objective while forcing:
% \begin{align}
% \vxx_{[end]} \cdot \ones_{ns} \leq \ME_\MZ \cdot \vxx  \iff [\zero \;\; \ones_{ns}] \vxx \leq \ME_\MZ \cdot \vxx
% \end{align}
% which enforces the extra variable to match the smallest entry of $\MZ$.


% ===================================================================
\section{Secondary AC Assignment Problem}

Now we tackle a related problem, where we want to assign a secondary AC to each paper. This is typically done \emph{before} creating the sessions as described above. We formulate the problem as follows:
\begin{align}
\max_{\substack{\MS \in \{0,1\}^{n \times p} \\ \MZ \in \{0,1\}^{n \times p} } } & \ones\tran \cdot \MZ \cdot \ones & \text{(maximize cumulative affinity score)} \nonumber \\
\subject & \MS \cdot \ones_p  \leq p_{\max,2} \cdot \ones_n\tran & \text{(at most $p_{\max,2}$ assignments per AC)} \nonumber \\
         &  \ones_n\tran \cdot \MS  = \ones_p  & \text{(exactly 1 secondary AC per paper)} \nonumber \\
         \nonumber \\
         &  \MS \leq \ones\tran \ones - \MP  & \text{(primary ACs cannot be secondary ACs for the same paper)} \nonumber \\
         &  \MS \leq \ones\tran \ones - \MC  & \text{(avoid COI)} \nonumber \\
         \nonumber \\ 
         & \MZ \leq a_\max \MS & \text{(force $\MZ$ to zero outside assignments)} \nonumber \\
         & \MZ \leq \MA & \text{(force upper-bound on $\MZ$ using the affinities)} \nonumber 
\end{align}
After reformulating the problem in terms of $\vxx = [\text{vec}(\MS) \; \text{vec}(\MZ)]$:
\begin{align}
\max_{\vxx \in \{0,1\}^{2 np}} & \ones_{np}\tran \cdot \ME_\MZ \cdot \vxx  \myQuad \text{(maximize cumulative affinity score)} \nonumber \\
\subject &  (\ones_p\tran \kron \eye_n) \cdot \ME_\MS \cdot \vxx \leq p_{\max,2} \cdot \ones_n & \text{(at most $p_{\max,2}$ assignments per AC)} \nonumber \\
         &  (\eye_p \kron \ones_n\tran)  \cdot \ME_\MS \cdot \vxx = \ones_p  & \text{(exactly 1 secondary AC per paper)} \nonumber \\
         \nonumber \\
         &  \ME_\MS \cdot \vxx \leq \text{vec}(\ones\tran \ones - \MP)  & \text{(primary ACs cannot be secondary ACs for the same paper)} \nonumber \\
         &  \ME_\MS \cdot \vxx \leq \text{vec}(\ones\tran \ones - \MC)  & \text{(avoid COI)} \nonumber \\
         \nonumber \\ 
         & \ME_\MZ \cdot \vxx \leq a_\max \ME_\MS \cdot \vxx & \text{(force $\MZ$ to zero outside assignments)} \nonumber \\
         & \ME_\MZ \cdot \vxx \leq \text{vec}(\MA) & \text{(force upper-bound on $\MZ$ using the affinities)} \nonumber 
\end{align}
where we used again the property of vectorization $\text{vec}(\MF \cdot \MG) =  (\MG\tran \kron \eye_a) \cdot \text{vec}(\MF) = 
(\eye_c \kron \MF) \cdot \text{vec}(\MG)$ for two matrices $\MF \in \Real{a \times b}$ and $\MG \in \Real{b \times c}$.
Moreover, we defined the selection matrices $\ME_\MS = [\eye_{np} \;\; \zero_{np,np}]$ and $\ME_\MZ = [\zero_{np} \;\; \eye_{np,np}]$ that select the entries corresponding to $\MS$ and $\MZ$ from $\vxx$, respectively.

% The decision variable in this problem is $\MX \in \{0,1\}^{n \times m}$, where $n$ is the number of Area Chairs (ACs) and $m$ is the number of sessions.
% In our case, $n = 54$ and $m = 33$ (4 paper discussion slots, each with 8 parallel breakouts + 1 ``standout'' session for ACs that cannot join any of the regular sessions).
% In particular $\MX_{ij} = 1$ if AC $i$ joins session $j$.
% %
% We are also given a matrix $\MA \in \Real{n \times m}$, which contains the affinity of each AC versus each session. In other words, 
% $\MA_{ij}$ is the affinity of AC $i$ for session $j$ (in practice this is computed as the average affinity over all the papers in the session, which should be thematically similar). The last column of $\MA$ is all zeros, meaning that there is no incentive for an AC to join the ``standout'' session. 
% Moreover, for the ACs that have a Conflict of Interest (COI) with at least a paper in that session, we set the corresponding entry of $\MA$ to $-\infty$ such that the optimization will never pick such an option.
% % 
% Therefore, the goal is to find the matching of ACs to sessions that maximizes the affinity scores, while satisfying constraints (\eg avoid COI, an AC can join one session at a time, etc.). 

% Next, we state the optimization problem that will compute the optimal matching:

% \begin{align}
% \max_{\MX \in \{0,1\}^{n \times m}} & \ones\tran (\MA \odot \MX) \ones & \text{(maximize cumulative affinity score)} \\
% \subject &  \MX \cdot \ve_{1:8} \leq \ones  & \text{(join only 1 of 8 breakouts for the 1st discussion slot)} \\
%          &  \MX \cdot \ve_{9:16} \leq \ones  & \text{(join only 1 of 8 breakouts for the 2nd discussion slot)} \\
%          &  \MX \cdot \ve_{16:24} \leq \ones  & \text{(join only 1 of 8 breakouts for the 3rd discussion slot)} \\
%          &  \MX \cdot \ve_{25:32} \leq \ones  & \text{(join only 1 of 8 breakouts for the 4th discussion slot)} \\
%          &  \ones\tran \MX = n \ones\tran & \text{(all ACs should join at least one session)} \\
% \end{align}
% where $\odot$ is the Hadamard (\ie element-wise) product, $\ones$ is the vector with all entries equal to one, 
% $\ve_{a:b}$ is the vector with entries from $a$ to $b$ equal to $1$ and zero otherwise.

% In Matlab's {\tt intlinprog}, the decision variable must be a vector, hence we rewrite the previous problem as:

% \begin{align}
% \max_{\vxx \in \{0,1\}^{nm}} & \va\tran \cdot \vxx  & \text{(maximize cumulative affinity score)} \\
% \subject &  (\ve_{1:8}\tran \kron \eye_n) \cdot \vxx \leq \ones  & \text{(join only 1 of 8 breakouts for the 1st discussion slot)} \\
%          &  (\ve_{9:16}\tran \kron \eye_n) \cdot \vxx \leq \ones  & \text{(join only 1 of 8 breakouts for the 2nd discussion slot)} \\
%          &  (\ve_{17:24}\tran \kron \eye_n) \cdot \vxx \leq \ones  & \text{(join only 1 of 8 breakouts for the 3rd discussion slot)} \\
%          &  (\ve_{25:32}\tran \kron \eye_n) \cdot \vxx \leq \ones  & \text{(join only 1 of 8 breakouts for the 4th discussion slot)} \\
%          &  (\ones\tran \kron \eye_m) \cdot \vxx = n \ones & \text{(all ACs should join at least one session)} \\
% \end{align}
% where $\va = \text{vec}(\MA)$ (column-wise vectorization), $\eye_n$ is the identity matrix of size $n$, and we used the property 
% that for two matrices $\MF$ and $\MG$ it holds $\text{vec}(\MF \cdot \MG) =  (\MG\tran \kron \eye_n) \cdot \text{vec}(\MF) = 
% (\eye_n \kron \MF) \cdot \text{vec}(\MG)$.


% assign papers to sessions

% AC and secondary AC for a paper must be in that session

% sessions must be balanced in terms of number of ACs

% \bibliographystyle{unsrt}
% \bibliography{/Users/lucacarlone/Desktop/myPapers/references/refs.bib} 

% that's all folks


\end{document}


